\chapter*{Before you start}
\addcontentsline{toc}{chapter}{Before you start}

\section*{How this software was made}
We developed \soft with the help of an artificial intelligence model, Claude by Anthropic, used as a programming
tool. We would rather say so clearly, once, here. The split was simple: we decided what the software should do, what
it should show and how to show it; we tested every build on our own hardware, with our own music, and turned down
whatever did not suit us. The tool wrote the code, the tests and the control renders from our instructions. No
build reaches you without being compiled, tested automatically, then tried on a real phone.

The work described in this manual runs from 2 October to 24 November 2026.

\section*{What you will find here}
This manual explains how to use \soft, but also what happens inside each sound chip and how to read it on screen.
If you only want to listen to your music, the first two parts are enough. If you compose, study the hardware or want
to understand why a tune sounds the way it does, the rest is for you.

A few conventions: registers are written in hexadecimal with the \$ prefix, notes in English notation (concert A is
A4, at 440~Hz) and levels in decibels relative to the chip's maximum. Screenshots are renders of the application on
a screen 1080 pixels wide.
